\documentclass[11pt,a4paper]{article}

% --- Packages ---
\usepackage[margin=.75in]{geometry}
\usepackage[utf8]{inputenc}
\usepackage[T1]{fontenc}
\usepackage{lmodern} % Fix for font not found error
\usepackage{titlesec}
\usepackage{enumitem}
\usepackage{hyperref}
\usepackage{xcolor}
\usepackage{fancyhdr}
\usepackage{comment}
\usepackage{cleveref}
\usepackage[bottom]{footmisc}
\usepackage{graphicx}
\usepackage{longtable}
\usepackage{booktabs}
\usepackage{array}
\usepackage{calc}
\usepackage{etoolbox}

% Standard Pandoc setup
\providecommand{\tightlist}{%
  \setlength{\itemsep}{0pt}\setlength{\parskip}{0pt}}

\crefformat{section}{\S#2#1#3}
\crefformat{subsection}{\S#2#1#3}
\crefformat{subsubsection}{\S#2#1#3}
\crefformat{figure}{#2Figure~#1#3}
\crefformat{footnote}{#2\footnotemark[#1]#3}

\linespread{0.9} % Riduce lo spazio tra le linee

% --- Colors ---
\definecolor{primary}{RGB}{0, 0, 0}
\definecolor{links}{HTML}{0000EE}
\definecolor{blue}{HTML}{26428B}

\hypersetup{
    colorlinks=true,
    linkcolor=links,
    urlcolor=links,
    citecolor=links,
}

% --- Section Formatting ---
\titleformat{\section}
  {\normalfont\Large\bfseries\color{blue}}{\thesection}{1em}{}
\titleformat{\subsection}
{\normalfont\large\bfseries\color{blue}}{\thesubsection}{1em}{}
\titleformat{\subsubsection}
{\normalfont\bfseries\color{blue}}{\thesubsubsection}{1em}{}

% --- Header and Footer ---
\pagestyle{fancy}
\fancyhead[L]{\small Giuseppe Capasso}
\fancyhead[R]{\small intro-to-linux}
\fancyfoot[C]{\thepage}
\renewcommand{\headrulewidth}{0.4pt}

% --- Custom Title Command ---
\newcommand{\proposaltitle}[1]{
    \begin{center}
        \vspace*{0.1cm}
        {\LARGE \bfseries \color{blue} #1} \\[1.5ex]
        {\large \textbf{Giuseppe Capasso}} \\[1ex]
                \small{
            \vspace{0.1cm}
            \textbf{Materia:} intro-to-linux\\
        }
            \end{center}
    \vspace{0.3cm}
    \hrule height 0.5pt
    \vspace{0.5cm}
}

\begin{document}

\proposaltitle{Scarichiamo la versione 3.3.11 (o latest stabile)}

Esercitazione Networking: ``The Double Identity'' Obiettivo: Configurare
indirizzi IP multipli, manipolare la risoluzione dei nomi (DNS locale) e
testare la connettività specifica. Scenario: Sei l'amministratore di un
server che deve offrire due servizi diversi: 1. Un servizio pubblico
(accesso Internet). 2. Un servizio ``Intranet Segreta'' accessibile solo
su una rete privata specifica. Non hai budget per una seconda scheda di
rete. Devi configurare la tua unica interfaccia per gestire una ``doppia
identità''.

Parte 1: Ricognizione (Discovery) Prima di toccare qualsiasi cosa, devi
capire la situazione attuale. 1. Trova il nome della tua interfaccia:
Usa il comando moderno (non ifconfig!) per elencare le interfacce. •
Comando: ip addr show (o ip a). • Nota: Cerca l'interfaccia che ha un IP
(solitamente enp0s3, eth0 o ens33). Ignora lo (loopback). 2. Identifica
il tuo IP attuale: Segnati l'indirizzo IP dinamico che VirtualBox ti ha
assegnato (probabilmente simile a 10.0.2.15). 3. Identifica il Gateway:
Scopri chi ti dà accesso a Internet. • Comando: ip route • Cerca la riga
che inizia con default via \ldots.

Parte 2: IP Aliasing (L'Identità Segreta) Ora aggiungeremo manualmente
un secondo indirizzo IP statico alla stessa scheda di rete. Questo IP
apparterrà a una rete completamente inventata (192.168.99.x). 1.
Aggiungi l'IP secondario: Sostituisci {[}NOME\_INTERFACCIA{]} con il
nome trovato nella Parte 1 (es. enp0s3). Bash sudo ip addr add
192.168.99.10/24 dev {[}NOME\_INTERFACCIA{]}

2. Verifica: Lancia di nuovo ip a. • Dovresti vedere due righe ``inet''
sotto la stessa interfaccia. • Una è quella vecchia (DHCP), l'altra è la
tua nuova 192.168.99.10. 3. Ping Test (Loopback): Prova a pingare te
stesso sul nuovo IP. Bash ping 192.168.99.10

(Funziona? Bene, il kernel ha accettato l'indirizzo).

Parte 3: DNS Spoofing Locale (/etc/hosts) I server non si chiamano per
numero, ma per nome. Faremo finta che questo nuovo IP corrisponda al
dominio segreto dell'azienda. 1. Apri il file hosts: Bash sudo nano
/etc/hosts

\begin{enumerate}
\def\labelenumi{\arabic{enumi}.}
\setcounter{enumi}{1}
\tightlist
\item
  Aggiungi in fondo al file questa riga: Plaintext 192.168.99.10
\end{enumerate}

portale-segreto.corp

\begin{enumerate}
\def\labelenumi{\arabic{enumi}.}
\setcounter{enumi}{2}
\tightlist
\item
  Salva (CTRL+O) ed esci (CTRL+X).
\item
  Test di risoluzione: Prova a pingare il nome invece del numero: Bash
  ping portale-segreto.corp
\end{enumerate}

(Se risponde l'IP 192.168.99.10, hai configurato correttamente il DNS
locale).

Parte 4: Simulazione Servizio Dedicato Ora avviamo un Web Server Python
che risponde SOLO sulla rete segreta, ignorando quella pubblica. 1. Crea
una cartella per il sito segreto: Bash mkdir /tmp/top\_secret echo ``

ACCESSO AUTORIZZATO AL PORTALE SEGRETO

'' \textgreater{} /tmp/top\_secret/index.html

2. Binding specifico: Avvia il server Python dicendogli di ascoltare
(bind) solo sull'IP 192.168.99.10 e non sull'IP pubblico. Bash cd
/tmp/top\_secret python3 -m http.server 8080 --bind 192.168.99.10

(Nota: Non chiudere questo terminale). 3. Il Test Finale: Apri un
secondo terminale (o usa un'altra tab). • Prova a scaricare la pagina
usando il nome DNS creato prima: Bash curl
http://portale-segreto.corp:8080

• Risultato atteso: Vedi il codice HTML ``ACCESSO AUTORIZZATO\ldots{}''.
• Prova ora a collegarti usando l'IP ``normale'' della VM (es.
10.0.2.15) sulla porta 8080: Bash curl http://10.0.2.15:8080

• Risultato atteso: Connection Refused o timeout. Perché? Perché abbiamo
configurato il servizio per ascoltare solo sull'identità segreta!

Parte 5: Pulizia (Reset) Poiché abbiamo usato il comando ip addr add
(che è volatile), per pulire tutto basta un comando o un riavvio. Ma
facciamolo da veri SysAdmin, manualmente. 1. Rimuovi l'IP secondario:
Bash sudo ip addr del 192.168.99.10/24 dev {[}NOME\_INTERFACCIA{]}

\begin{enumerate}
\def\labelenumi{\arabic{enumi}.}
\setcounter{enumi}{1}
\tightlist
\item
  Rimuovi la riga da /etc/hosts (con nano).
\item
  Verifica con ip a che tutto sia tornato come prima.
\end{enumerate}

Mini-Progetto: Deploy di un Forum Aziendale (phpBB) Scenario: Il team di
supporto tecnico ha bisogno di una piattaforma interna per discutere dei
ticket e condividere soluzioni (Knowledge Base). Ti è stato chiesto di
installare un forum phpBB sul server di test. Devi configurare il server
LAMP, il database e i permessi del filesystem affinché l'applicazione
funzioni correttamente. Obiettivo: Avere un forum funzionante
accessibile via browser all'indirizzo http://{[}IP\_VM{]}/forum. Tempo
stimato: 60-70 Minuti

Fase 1: Preparazione dello Stack (Dipendenze) phpBB è un software
robusto che richiede moduli PHP specifici per gestire immagini (avatar)
e file XML. 1. Aggiorna e installa LAMP di base: Bash sudo apt update
sudo apt install apache2 mariadb-server php libapache2-mod-php php-mysql
unzip wget -y

\begin{enumerate}
\def\labelenumi{\arabic{enumi}.}
\setcounter{enumi}{1}
\item
  Installa i moduli PHP specifici per phpBB: Senza questi,
  l'installazione si bloccherà. Bash sudo apt install php-xml
  php-mbstring php-gd php-intl php-curl php-zip -y
\item
  Riavvia Apache per caricare i nuovi moduli: Bash sudo systemctl
  restart apache2
\end{enumerate}

Fase 2: Configurazione del Database 1. Accedi a MariaDB: Bash sudo mysql

\begin{enumerate}
\def\labelenumi{\arabic{enumi}.}
\setcounter{enumi}{1}
\tightlist
\item
  Crea il database e l'utente dedicato (Esegui una riga alla volta): SQL
\end{enumerate}

CREATE DATABASE forum\_db; CREATE USER `forum\_admin'@`localhost'
IDENTIFIED BY `password\_sicura'; GRANT ALL PRIVILEGES ON forum\_db.* TO
`forum\_admin'@`localhost'; FLUSH PRIVILEGES; EXIT;

Fase 3: Download e Setup Filesystem A differenza dell'esercizio
precedente, qui non metteremo il sito nella root (/), ma in una
sottocartella /forum. 1. Vai in /tmp e scarica l'ultima versione di
phpBB (in italiano): Bash cd /tmp \# Scarichiamo la versione 3.3.11 (o
latest stabile) wget
https://download.phpbb.com/pub/release/3.3/3.3.11/phpBB-3.3.11.zip

\begin{enumerate}
\def\labelenumi{\arabic{enumi}.}
\setcounter{enumi}{1}
\tightlist
\item
  Scompatta: Bash unzip phpBB-3.3.11.zip
\end{enumerate}

(Questo creerà una cartella chiamata phpBB3). 3. Deploy: Sposta la
cartella nella directory web e rinominala in ``forum'': Bash sudo mv
phpBB3 /var/www/html/forum

Fase 4: Gestione Avanzata dei Permessi phpBB è molto severo sulla
sicurezza. Alcune cartelle devono essere scrivibili dal server web
(wwwdata), altre no. 1. Assegna la proprietà: Tutto appartiene a
www-data (l'utente di Apache). Bash sudo chown -R www-data:www-data
/var/www/html/forum

\begin{enumerate}
\def\labelenumi{\arabic{enumi}.}
\setcounter{enumi}{1}
\tightlist
\item
  Permessi specifici (Security Hardening): Di base, i file devono essere
  leggibili ma non modificabili. Le cartelle di cache e upload devono
  essere scrivibili. Imposta permessi standard (755 per cartelle, 644
  per file):
\end{enumerate}

Bash sudo find /var/www/html/forum -type d -exec chmod 755 \{\} ; sudo
find /var/www/html/forum -type f -exec chmod 644 \{\} ;

Apri i permessi solo dove serve (Cache, Store, Files, Images): Bash sudo
chmod 777 /var/www/html/forum/cache sudo chmod 777
/var/www/html/forum/store sudo chmod 777 /var/www/html/forum/files sudo
chmod 777 /var/www/html/forum/images/avatars/upload

(Nota: In produzione useremmo permessi più stretti tipo 775 con il
gruppo corretto, ma per questo lab 777 su queste specifiche cartelle
interne va bene per capire il concetto di ``writable'').

Fase 5: Installazione via Browser 1. Apri il browser sul tuo PC Host. 2.
Digita: http://{[}IP\_DELLA\_TUA\_VM{]}/forum 3. Dovresti vedere la
pagina di installazione di phpBB. Clicca sulla scheda Installazione
(Install). 4. Segui il wizard: • Dati Amministratore: Crea il tuo utente
(es. admin, pass: admin123). • Database: • Tipo: MySQL with MySQLi
Extension • Host: localhost • Nome DB: forum\_db • Utente: forum\_admin
• Password: password\_sicura • Configurazione Server: Lascia i default.
5. Prosegui fino alla fine dell'installazione.

Fase 6: Pulizia Obbligatoria (Security) Alla fine dell'installazione,
phpBB ti mostrerà un avviso rosso gigante: ``Rimuovere, spostare o
rinominare la cartella install''. Finché non lo fai, il forum sarà
bloccato per sicurezza. 1. Torna sul terminale della VM. 2. Elimina la
cartella di installazione: Bash sudo rm -rf /var/www/html/forum/install

\begin{enumerate}
\def\labelenumi{\arabic{enumi}.}
\setcounter{enumi}{2}
\tightlist
\item
  Torna sul browser e ricarica la pagina. Il tuo forum è online!
\end{enumerate}

🔎 Troubleshooting • Errore 500 o pagina bianca: Controlla se hai
installato php-xml e riavviato Apache. È l'errore più comune con phpBB.
• ``Directory is not writable'': Hai sbagliato la Fase 4. Ricontrolla i
comandi chmod. • Impossibile connettersi al DB: Hai usato root invece
dell'utente forum\_admin? In alcune installazioni di MariaDB, root non
può accedere via password esterna per sicurezza. Usa sempre l'utente
dedicato.

★ Challenge Extra (Per i più veloci) Customizzazione del messaggio di
benvenuto (HTML/CSS Base) Il forum funziona, ma è brutto. 1. Entra nel
Pannello di Controllo Amministrazione (Link in fondo alla pagina del
forum -\textgreater{} ``ACP''). 2. Vai su ``Generale'' -\textgreater{}
``Configurazione Board''. 3. Cambia il nome del sito e la descrizione.
4. Livello Pro: Prova a sostituire il logo di phpBB via terminale. •
Trova un'immagine .png su internet. • Scaricala con wget dentro
/var/www/html/forum/styles/prosilver/theme/images/. • Rinominala per
sovrascrivere site\_logo.svg (o modifica il CSS se sei capace).



\end{document}
